\renewcommand{\thesection}{A-\Roman{section}}
\renewcommand{\thetable}{A-\Roman{table}}
\renewcommand{\thefigure}{A-\arabic{figure}}
\setcounter{section}{0}
\setcounter{table}{0}
\setcounter{figure}{0}


\begin{foldable} % Appendix overview
    In this extended piece, we provide our implementation details, complementary studies on compression ratio and alternative optimization objectives for our image priors, and extension to an iterative framework.
\end{foldable}


\section{Training Details} \label{sec:app01-training}
\subsection{Datasets} \label{ssec:app01-training-datasets}
\begin{foldable} % Datasets details
    In our experiments (Sec.~\ref{sec:04-experiments}), we evaluate our method on eight benchmark image datasets: MNIST~\cite{lecun2002gradient}, USPS~\cite{hull1994database}, SVHN~\cite{netzer2011reading}, FMNIST~\cite{xiao2017fashion}, CIFAR10, CIFAR100~\cite{krizhevsky2009learning}, TinyImgN~\cite{le2015tiny}, Imagenette~\cite{howard2020imagenette}.
    All datasets are public and easily downloable with standard libraries.
    In terms of content, MNIST, USPS, and SVHN contain small images of digits.
    FMNIST is also a small-sized dataset of fashion items.
    The images in CIFAR10, CIFAR100, TinyImgN, and Imagenette all contain diverse scenes and objects, e.g., animals, plants, food, vehicles, and household goods, to name a few.
    Notably, TinyImgN and Imagenette have higher resolution: 64$\times$64 for TinyImgN, and approximately one magnitude larger than that for Imagenette (ranging from 27$\times$80 to 4368$\times$2912 --- which we reported the average dimension).
    Of all benchmarks, only CIFAR100 and TinyImgN have abundant classes (100 and 200, respectively), while the others have 10 classes.
    Regarding the dataset size, USPS and Imagenette are relatively smaller.
    We summarize their details in Tab.~\ref{tab:appendix-training-datasets}.
\end{foldable}

\begin{table}[ht] % tab:appendix-training-datasets
    \caption{Details of the Image Benchmark Datasets}
    \centering
    \begin{tabular}{c|c|c|c|c}
        \hline \textbf{Dataset} & \textbf{Dimension} & \textbf{Size} ($\times\text{10}^\text{3}$) & \textbf{No.} & \textbf{Content}  \\
        & & train; test & \textbf{classes} & \\
        \hline
        % \hline
        MNIST      &                    1$\times$28$\times$28 &           \phantom{0}60 ; 10\phantom{0} & \phantom{0}10 & Digits \\
        USPS       &                    1$\times$16$\times$16 & \phantom{00}9 ; \phantom{0}2\phantom{0} & \phantom{0}10 & Digits \\
        SVHN       &                    3$\times$32$\times$32 &           \phantom{0}73 ; 26\phantom{0} & \phantom{0}10 & Digits \\
        FMNIST     &                    1$\times$28$\times$28 &           \phantom{0}60 ; 10\phantom{0} & \phantom{0}10 & Fashion items \\
        \hline
        % \multicolumn{5}{c}{\textbf{Complex datasets}} \\
        % \hline
        CIFAR10    &                    3$\times$32$\times$32 &           \phantom{0}50 ; 10\phantom{0} & \phantom{0}10 & Diverse \\
        CIFAR100   &                    3$\times$32$\times$32 &           \phantom{0}50 ; 10\phantom{0} &           100 & Diverse \\
        TinyImgN   &                    3$\times$64$\times$64 &                     100 ; 10\phantom{0} &           200 & Diverse \\
        Imagenette & 3$\times$408$\times$472$^\blacktriangle$ & \phantom{00}9 ; \phantom{0}4\phantom{0} & \phantom{0}10 & Diverse \\
        \hline
        \multicolumn{5}{l}{$^\blacktriangle$: The average dimension is reported.} \\
    \end{tabular}
    \label{tab:appendix-training-datasets}
\end{table}


\subsection{Teacher and students training} \label{ssec:app01-teacherstudents}
\subsubsection{Augmentation} \label{sssec:app01-teacherstudents-augment}
\begin{foldable} % Augmentation
    To augment the training data (real images for the teacher, synthetic images for students), we resize very small images (in MNIST, USPS, FMNIST) to at least $32\times32$.
    Then the images undergo random shifting by 1/8 of their dimension, and random horizontal flipping unless they are from a digits dataset (MNIST, USPS, SVHN).
    Finally, we randomly jitter the color in the images for more image variety.
    For the high-resolution Imagenette, we efficiently perform $4\times$ downscale at storage and $4\times$ upscale at runtime.
\end{foldable}

\subsubsection{Optimization and scheduling} \label{sssec:app01-teacherstudents-optimization}
\begin{foldable} % Optimization and scheduling
    To train the Teacher for baseline, and the Students in our method, we use a stochastic gradient descent optimizer with momentum of 0.9, weight decay of $5\times10^{-4}$, and batch size of 100.
    On datasets with small black-and-white images (MNIST, USPS, FMNIST), we use a fixed learning rate of $10^{-2}$ and train for 100 epochs.
    On other datasets, we set the learning rate to $10^{-1}$ at the beginning and decrease by two orders of magnitude with a scheduler to enable optimal convergence for 200 training epochs.
\end{foldable}


\subsection{Optimization of Image Priors} \label{ssec:app01-imagepriors}
\subsubsection{Diverging filter} \label{sssec:app01-teacherstudents-optimization-divgfilter}
\begin{foldable} % Diverging filter principles
    We provide more details on the formulation of the diverging filter $f_{\text{divg}}$ (Eq.~\ref{eq:synthesis-diverging-filter}) proposed in the Synthesis phase (Sec.~\ref{ssec:03-framework-synthesis}).
    To fulfill its role in contrasting opposite-type regions in a mask $m\in[0, 1]^{C\times H\times W}$, we have the following set of principles:
    \begin{enumerate}
        \item We start by picking a brightness cutoff threshold $\beta$ to separate the negative ($x<\beta$) and postive regions ($x\ge\beta$) for any pixel $x\in m$.
            The extrema and the threshold value must be fixed after filtering, \ie, $f_{\text{divg}}\left(\cdot; \beta\right)$ must go through $\left\{ (0,0), (\beta,\beta), (1,1)\right\}, \forall\beta$.
        \item To preserve the relative intensity across the whole mask, it is necessary that $f_{\text{divg}}$ is an increasing function, \ie, $f_{\text{divg}}^{\prime}\left(x;\beta\right)>0,\forall x,\beta\in\left[0,1\right]$.
        \item To avoid creating abrupt changes in intensity, \ie, $f_{\text{divg}}$ must be continuously differentiable on $[0,1], \forall\beta$.
        \item To ensure the diverging property emerges at the threshold value $x=\beta$, we require that $f_{\text{divg}}$ has the sharpest increase there, \ie, $f_{\text{divg}}^{\prime}(\cdot;\beta)$ has a maximum at $x=\beta\Leftrightarrow f^{\prime\prime}(x=\beta;\beta)=0,f^{\prime}(x=\beta;\beta)>0$.
    \end{enumerate}
\end{foldable}

\begin{figure}[ht] % fig:diverging-filter
    \centering
    \includegraphics[width=0.7\columnwidth]{./figures/diverging-filter.pdf}
    \caption{
        Our diverging filter as $y=f_{\text{divg}}\left(x;\beta\right)$ in Eq.~\ref{eq:synthesis-diverging-filter} with $x,\beta\in\left[0,1\right]$.
        Scattered points $\left\{\left(\beta,\beta\right)\right\}$ along $y=x$ are also inflection points.
    }
    \label{fig:diverging-filter}
\end{figure}

\begin{foldable} % Diverging filter design
    Accordingly, we elegant design a piecewise function consisting of two continuously differentiable half-parabolas connected at the inflection point $\left(\beta,\beta\right)$.
    One of them has a minimum at the $x=0$ boundary, and vice versa, the other half-parabola has a maximum at the opposite boundary $x=1$, \ie, $f_{\text{divg}}^{\prime}(0;\beta)=f_{\text{divg}}^{\prime}(1;\beta)=0$.
    Solving the two parabolas with the proposed set of constraints indeed leads us back to the solution we reported in Eq.~\ref{eq:synthesis-diverging-filter}.
    We plot $f_{\text{divg}}$ with a range of $\beta$ values in Fig.~\ref{fig:diverging-filter} to demonstrate its diverging property.
    We believe that this filter design has unexplored potentials in image augmentation for robust learning of vision tasks.
\end{foldable}



\subsubsection{Loading unique embeddings during Contrastion} \label{sssec:app01-teacherstudents-optimization-uniqueembeddings}
\begin{foldable} % Problem: exploding gradients
    Originally, the contrastive method in~\cite{fang2021contrastive} store the embeddings of past inputs $x$ and their corresponding positive views $x^{+}$ in a memory bank, while the negative views $x^{-}$ are obtained by pairing different images together.
    However, on with small datasets, we occasionally observe numerical instability when an embedding loaded from the memory bank is near-duplicated with one extracted at runtime (different by some minor updates) as they may originate from the same image.
    Optimization-wise, contrasting an image against itself may lead to exploding gradients and corrupt the synthetic images.
\end{foldable}

\begin{foldable} % Alternative fix
    Therefore, we modify how images are loaded for contrastive optimization.
    Instead of using a memory bank, we store the embeddings of each synthetic images in our dataset by its unique index.
    At runtime, we only load embeddings of images \textbf{not} in the current batch to perform contrastion.
    This modification fixes the numerical instability on our ends.
\end{foldable}


\subsubsection{Instance-discriminator} \label{sssec:app01-teacherstudents-optimization-instancedisc}
\begin{foldable} % Instance-discriminator
    Following~\cite{fang2021contrastive}, we use a fully connected instance-discriminator $R$.
    For each input image $x$, $R$ projects the image embedding $h=S^{\mathcal{H}}(x)$ from student backbone $S^{\mathcal{H}}$, through one hidden layer, to its own embedding space of 128 dimensions.
    To update $R$ during the Contrastion phase (Sec.~\ref{ssec:03-framework-contrastion}), we backpropagate the contrastive loss $\mathcal{L}_{\text{CT}}$ in Eq.~\ref{eq:loss-contrastive} to it parameters.
    It is trained with an Adam optimizer~\cite{kingma2014adam} with LR of $10^{-3}$, and its ``betas'' coefficients of (0.5, 0.999).
    As $R$ is shallow (two linear layers), this does not incur a significant computational overhead.
\end{foldable}

\subsubsection{Resources} \label{sssec:app01-teacherstudents-optimization-resources}
We perform each standard experiment on a single NVIDIA V100 GPU with 40GB VRAM, with at most 16 CPUs, and 4--128 GB RAM.
Each run takes around 6--96 hours.
The training time is proportional to the model size, dataset size, and synthetic image resolution.


\section{Complementary Studies} \label{sec:app02-experiments}
\begin{foldable} % Complementary studies overview
    In the main paper (Sec.~\ref{sec:04-experiments}), we have investigated BBDFKD with standard and extended settings (proxy data, cross-architecture), inclusive ablation study, varying synthetic dataset size, and embedding analysis.
    In the following section, we will examine student by varying its compression ratio and different objectives for image prior optimization.
\end{foldable}


\subsection{Students of compressed size} \label{ssec:app02-experiments-ablation-compress}
\begin{table}[ht]
    \caption{Parameters Count by Compression Ratio (CR)}
    \label{tab:abla-compression-parameters}
    \centering
    \begin{tabular}{c|c|c|c}
        \hline \textbf{CR} & \textbf{LeNet5} & \textbf{AlexNet} & \textbf{ResNet18} \\
        \hline        100\% &           61.71 K &           1,659.18 K &           11.17 M \\
            \phantom{0}50\% &           15.74 K & \phantom{0,}417.43 K & \phantom{0}2.90 M \\
            \phantom{0}20\% & \phantom{0}2.53 K & \phantom{0,0}68.49 K & \phantom{0}0.45 M \\
            \phantom{0}10\% & \phantom{0}0.88 K & \phantom{0,0}17.70 K & \phantom{0}0.11 M \\
        \hline
    \end{tabular}
\end{table}

\begin{figure}[ht] % fig:abla-compression
    \centering
    \includegraphics[width=\columnwidth]{./figures/experiments-abla-compression.pdf}
    \caption{Distillation accuracy (\%) by compression ratio on simple datasets.}
    \label{fig:abla-compression}
\end{figure}

\begin{foldable}
    We investigate the case where the student architecture gets compressed.
    This is done by reducing the nodes in linear layers or the channels in convolutional layers by a ratio of 50\%, 20\%, 10\%.
    Note that the actual parameter count in the model decreases near-quadratically by these ratios (Tab.~\ref{tab:abla-compression-parameters}).
    We report the KD results in Fig.~\ref{fig:abla-compression}.
    Observably, compression corrupts student performance, where LeNet5 and ResNet18 diminish at around 20\% compression ratio, and AlexNet still holds it until 10\% compression ratio before a large drop.
    However, performance of our 20\%-compressed students (0.04$\times$ parameters) are often acceptable. 
\end{foldable}


\subsection{Multi-objective image optimization} \label{ssec:app02-experiments-multiobjective}
\begin{foldable} % In-class similarity: One-hot loss
    Unlike the objective of diversity in our framework, previous BBDFKD methods such as~\cite{zhang2022ideal, yuan2024data} attempted to promote in-class similarity and class-balance.
    Specifically, in-class similarity is attained via the one-hot loss:
    \begin{equation} \label{loss-onehot}
        \mathcal{L}_{\text{OH}}=\mathcal{L_{\text{CE}}}\left(\hat{y}_{S-\text{Soft}},\text{argmax}(\hat{y}_{S-\text{Soft}})\right),
    \end{equation}
    which guides data generation to strengthen the effect of the most prominent class.
    The one-hot loss is popular in data-free KD~\cite{chen2019data}, but as stated in~\cite{yuan2024data}, its adversarial nature might cause ``overfitting of synthetic data to the clone model'' ($S_{0}$), resulting in mode collapse and insufficient diversity.
\end{foldable}

\begin{foldable} % Class-balance: Information entropy loss
    The other objective, class-balance, is to ensure the samples are evenly distributed between all classes, enforced by the information entropy loss:
    \begin{equation} \label{loss-infoentropy}
        \mathcal{L}_{\text{IE}}=\sum_{k=0}^{K-1}p_{k}\log p_{k},\text{ with }p_{k}=\mathbb{E}_{x\sim \mathcal{D}}\left[\hat{y}_{S-\text{Soft}}^{\left(k\right)}\right],
    \end{equation}
    where $k$ denotes the $k$-th class probability.
    In our framework, as we have actively balanced the dataset during sampling, class-balance has already been satisfied.
\end{foldable}

\begin{figure}[ht] % fig:abla-losses
    \centering
    \begin{minipage}{0.48\columnwidth}
        \centering
        \includegraphics[width=\columnwidth]{./figures/experiments-abla-losses-mnist-lenet5.pdf}
        {\footnotesize \textbf{(a)} MNIST}
    \end{minipage}
    \hfill
    \begin{minipage}{0.48\columnwidth}
        \centering
        \includegraphics[width=\columnwidth]{./figures/experiments-abla-losses-cifar10-resnet18.pdf}
        {\footnotesize \textbf{(b)} CIFAR10}
    \end{minipage}
    \caption{Contribution to distillation accuracy (\%) of one-hot loss ($\mathcal{L}_{\text{OH}}$), information entropy loss ($\mathcal{L}_{\text{IE}}$), and contrastive loss ($\mathcal{L}_{\text{CT}}$), with 50K samples.}
    \label{fig:abla-losses}
\end{figure}

\begin{foldable} % Objectives discussion
    To measure how these objectives might affect KD performance, we ran all eight combinations of using/not using $\mathcal{L}_{\text{OH}}$, $\mathcal{L}_{\text{IE}}$, $\mathcal{L}_{\text{CT}}$ (Fig.~\ref{fig:abla-losses}).
    The contribution of each objective is derived by the mean difference in student accuracy between using/not using the corresponding loss, while standard error is aggregated on a per-combination basis.
    As class-balance has already been achieved, $\mathcal{L}_{\text{IE}}$ only has modest contribution.
    Surprisedly, in-class similarity $\mathcal{L}_{\text{OH}}$ harms KD performance, significantly by $-$2.65\% on CIFAR10.
    Meanwhile, contrastion fruitfully helped gain KD performance, notably by $+$0.12\% on MNIST and $+$3.14\% on CIFAR10.
    Thus, we employ only the contrastive loss $\mathcal{L}_{\text{CT}}$ to optimize image priors, which only promotes diversity and is not adversarially destructive.
\end{foldable}

\begin{foldable} % Outro
    On a side note, \cite{zhang2022ideal, yuan2024data} employed $\mathcal{L}_{\text{OH}}$ and $\mathcal{L}_{\text{IE}}$ on the weights of a generator network, while we employ $\mathcal{L}_{\text{CT}}$ directly on the parameterized image priors $\mathcal{X}_{\mathcal{P}}$.
    We speculate that this difference may require further investigation.
\end{foldable}


\begin{figure}[ht] % fig:abla-ngenerations
    \centering
    \includegraphics[width=\columnwidth]{./figures/experiments-abla-ngenerations.pdf}
    \caption{Distillation accuracy (\%) by generation $g$.}
    \label{fig:abla-ngenerations}
\end{figure}


\section{Extension to Iterative Framework} \label{sec:app03-extension}
\begin{foldable}
    An interesting extension of our framework is performing \textit{iterative Contrastion-Distillation} to further boost student performance.
    Recall that in the Contrastion phase (Sec.~\ref{ssec:03-framework-contrastion}), we employ an initial student $S_0$ to optimize the dataset $\mathcal{D}$, and train a new student $S$ in the next Distillation phase.
    It naturally emerges that the stronger $S$ can be recycled for Contrastion again to produce even more robust embeddings.
    Iteratively, we may re-employ the latest student $S_g$ at the $g$\nobreakdash-th generation to optimize $\mathcal{D}$, and use it to train the new student $S_{g+1}$.
    The objectives for $S_{g+1}$, extended from $\mathcal{L}_S$ (Eq.~\ref{eq:loss-kd-total}), is to match the teacher's label $T(x)$ with $\mathcal{L}_\text{Hard-KD}$, and to match previous student's logits $S_{g}^{\mathcal{Z}}(x)$ with $\mathcal{L}_\text{Soft-KD}$, on samples $x\sim\mathcal{D}$.
\end{foldable}

\begin{foldable}
    From our standard experiments (Sec.~\ref{ssec:04-experiments-standard}, corresponds to $g$ = 2), we iterate for at most three more generations.
    As in Fig.~\ref{fig:abla-ngenerations}, KD performance converges very quickly on simple datasets, but still increases steadily $g$ = 3 with complex datasets.
    These results conclude that our iterative contrastion-distillation strategy can robustly deliver increasingly better students with more generations.
    Nonetheless, computational cost increases linearly but student accuracy does not scale as fast.
    We look forward to making this extension even more efficient in future works.
\end{foldable}


\begin{comment}


\section{Explicit Results} \label{sec:app04-explicit}
\subsection{Students of compressed size} \label{ssec:app04-explicit-compression}
\textbf{To rewrite due to changing datasets grouping.}
We provide in Tab~\ref{tab:abla-compression} explicit student accuracy and standard error for KD by compression ratio, which has been visualized and discussed in Fig.~\ref{fig:abla-compression}, Sec.~\ref{ssec:app02-experiments-ablation-compress}.

\begin{table}[ht] % tab:abla-compression
    \caption{Distillation accuracy (\%) by compression ratio (CR) \\ on Simple datasets \& CIFAR10}
    \centering
    \begin{tabular}{c|c|c|c|c}
        \hline
            {\centering \textbf{CR}}
            & \makecell{\textbf{MNIST}   \\ LeNet5}
            % & \makecell{\textbf{USPS}    \\ LeNet5}
            & \makecell{\textbf{SVHN}    \\ AlexNet}
            & \makecell{\textbf{FMNIST}  \\ LeNet5}
            & \makecell{\textbf{CIFAR10} \\ ResNet18} \\
        \hline
                  100\% & 98.59$_{\pm \phantom{0}.08}$ & 95.94$_{\pm \phantom{0}.05}$ & 83.98$_{\pm \phantom{0}.65}$ & 83.41$_{\pm \phantom{0}.46}$ \\
        \phantom{0}50\% &           96.79$_{\pm 1.39}$ & 95.41$_{\pm \phantom{0}.04}$ & 74.90$_{\pm \phantom{0}.37}$ & 79.23$_{\pm \phantom{0}.51}$ \\
        \phantom{0}20\% &           85.83$_{\pm 1.86}$ & 92.51$_{\pm \phantom{0}.07}$ &           52.55$_{\pm 3.79}$ & 64.48$_{\pm \phantom{0}.81}$ \\
        \phantom{0}10\% &           47.72$_{\pm 7.44}$ &           72.00$_{\pm 2.89}$ &           36.85$_{\pm 8.66}$ &           47.39$_{\pm 3.30}$ \\
        \hline
    \end{tabular}
    \label{tab:abla-compression}
\end{table}

\subsection{Number of samples. \textbf{To rewrite due to changing datasets grouping.}} \label{ssec:app04-explicit-nsamples}
\begin{foldable}
    We provide explicit student accuracy and standard error for KD by number of samples, which has been visualized and discussed in Fig.~\ref{fig:abla-nsamples}, Sec.~\ref{sssec:04-experiments-ablation-nsamples}.
    The results are grouped by the simple datasets in Tab.~\ref{tab:abla-nsamples-simple} and complex datasets in Tab~.\ref{tab:abla-nsamples-complex}.
\end{foldable}

\begin{table}[ht] % tab:abla-compression
    \caption{Distillation accuracy (\%) by number of sample $N$ \\ on simple datasets.}
    \centering
    \begin{tabular}{c|c|c|c|c}
        \hline
            $N$
            & \makecell{\textbf{MNIST}   \\ LeNet5}
            & \makecell{\textbf{USPS}    \\ LeNet5}
            & \makecell{\textbf{SVHN}    \\ AlexNet}
            & \makecell{\textbf{FMNIST}  \\ LeNet5} \\
        \hline
        \phantom{0}10 K &          98.20$_{\pm .06}$ &          93.22$_{\pm .43}$ &          94.64$_{\pm .13}$ &                    76.14$_{\pm 2.48}$ \\
        \phantom{0}20 K &          98.59$_{\pm .08}$ &          93.99$_{\pm .24}$ &          95.46$_{\pm .07}$ &                    79.49$_{\pm 1.38}$ \\
        \phantom{0}30 K &          98.68$_{\pm .03}$ &          94.05$_{\pm .08}$ &          95.71$_{\pm .05}$ &          82.67$_{\pm \phantom{0}.45}$ \\
        \phantom{0}50 K &          98.75$_{\pm .03}$ &          94.20$_{\pm .21}$ &          95.94$_{\pm .05}$ &          83.98$_{\pm \phantom{0}.65}$ \\
        \phantom{0}70 K &          98.81$_{\pm .03}$ &          94.25$_{\pm .06}$ &          95.97$_{\pm .02}$ &          84.84$_{\pm \phantom{0}.69}$ \\
                  100 K &          98.78$_{\pm .05}$ &          94.41$_{\pm .08}$ &          96.06$_{\pm .03}$ & \textbf{85.63}$_{\pm \phantom{0}.39}$ \\
                  128 K & \textbf{98.90}$_{\pm .02}$ & \textbf{94.44}$_{\pm .09}$ & \textbf{96.13}$_{\pm .01}$ &          85.62$_{\pm \phantom{0}.06}$ \\
        \hline
        \multicolumn{5}{l}{$^{(*)}$: Best results are in \textbf{bold}.} \\
    \end{tabular}
    \label{tab:abla-nsamples-simple}
\end{table}

\begin{table}[ht] % tab:abla-compression
    \caption{Distillation accuracy (\%) by number of samples $N$ \\ on complex datasets.}
    \centering
    \begin{tabular}{c|c|c|c|c}
        \hline
            $N$ & \textbf{CIFAR10} & \textbf{CIFAR100} & \textbf{TinyImgN} & \textbf{Imagenette} \\
        \hline
        \phantom{0}10 K &                    63.54$_{\pm 2.64}$ &                    25.85$_{\pm 1.31}$ &          07.88$_{\pm .49}$ &                    50.82$_{\pm 1.65}$ \\
        \phantom{0}20 K &          72.98$_{\pm \phantom{0}.50}$ &          37.48$_{\pm \phantom{0}.48}$ &          12.13$_{\pm .17}$ &                    57.43$_{\pm 1.51}$ \\
        \phantom{0}30 K &          78.18$_{\pm \phantom{0}.60}$ &          42.82$_{\pm \phantom{0}.56}$ &          13.96$_{\pm .42}$ &                    58.98$_{\pm 1.19}$ \\
        \phantom{0}50 K &          83.41$_{\pm \phantom{0}.46}$ &          48.45$_{\pm \phantom{0}.26}$ &          16.49$_{\pm .37}$ &          61.84$_{\pm \phantom{0}.77}$ \\
        \phantom{0}70 K &          85.20$_{\pm \phantom{0}.35}$ &          50.75$_{\pm \phantom{0}.18}$ &          18.69$_{\pm .63}$ &          63.08$_{\pm \phantom{0}.96}$ \\
                  100 K &          86.60$_{\pm \phantom{0}.29}$ &          52.85$_{\pm \phantom{0}.30}$ &          19.99$_{\pm .36}$ &          65.98$_{\pm \phantom{0}.78}$ \\
                  128 K &          87.56$_{\pm \phantom{0}.11}$ &          53.69$_{\pm \phantom{0}.26}$ &          21.07$_{\pm .50}$ & \textbf{66.60}$_{\pm \phantom{0}.72}$ \\
                  256 K & \textbf{89.24}$_{\pm \phantom{0}.15}$ & \textbf{56.84}$_{\pm \phantom{0}.54}$ &          24.46$_{\pm .26}$ &                                     - \\
                  512 K &                                     - &                                     - & \textbf{28.06}$_{\pm .22}$ &                                     - \\
        \hline
        \multicolumn{5}{l}{Best results are in \textbf{bold}.} \\
    \end{tabular}
    \label{tab:abla-nsamples-complex}
\end{table}


\subsection{Number of generations. \textbf{To rewrite due to changing datasets grouping.}} \label{ssec:app03-extension-ngenerations}
\begin{foldable}
    We provide explicit student accuracy and standard error for KD by number of generations, which has been visualized and discussed in Fig.~\ref{fig:abla-ngenerations}, Sec.~\ref{ssec:app03-extension-ngenerations}.
    We ran the experiments for the simple datasets, as shown in Tab.~\ref{tab:abla-ngenerations}.
    We additionally remark that the student can get iteratively better until the last generations, with some trivial fluctuations.
\end{foldable}

\begin{table*}[t] % tab:abla-generation
    \caption{Distillation accuracy (\%) by generation $g$}
    \centering
    \begin{tabular}{c|c|c|c|c|c|c|c|c}
        \hline
            $g$
            & \makecell{\textbf{MNIST}   \\ LeNet5}
            & \makecell{\textbf{USPS}    \\ LeNet5}
            & \makecell{\textbf{SVHN}    \\ AlexNet}
            & \makecell{\textbf{FMNIST}  \\ LeNet5}
            & \makecell{\textbf{CIFAR10} \\ ResNet18}
            & \makecell{\textbf{CIFAR100} \\ ResNet18}
            & \makecell{\textbf{TinyImgN} \\ ResNet18}
            & \makecell{\textbf{Imagenette} \\ ResNet18} \\
            \hline
        0 &          98.29$_{\pm .32}$ &            94.04$_{\pm .52}$ &          95.84$_{\pm .10}$ &                    82.38$_{\pm 1.09}$ &          78.17$_{\pm .64}$ &          48.36$_{\pm .19}$ &          22.66$_{\pm .36}$ &          57.67$_{\pm .74}$ \\
        1 &          98.59$_{\pm .08}$ &            94.20$_{\pm .21}$ &          95.94$_{\pm .05}$ &          83.98$_{\pm \phantom{0}.65}$ &          83.41$_{\pm .46}$ &          52.85$_{\pm .30}$ &          28.03$_{\pm .22}$ &          61.84$_{\pm .77}$ \\
        2 &          98.62$_{\pm .12}$ &            94.32$_{\pm .12}$ &          95.95$_{\pm .03}$ &          84.43$_{\pm \phantom{0}.66}$ &          84.50$_{\pm .16}$ & \textbf{53.97}$_{\pm .62}$ & \textbf{30.06}$_{\pm .22}$ & \textbf{62.78}$_{\pm .91}$ \\
        3 &          98.61$_{\pm .08}$ &            94.42$_{\pm .25}$ & \textbf{95.97}$_{\pm .06}$ &          84.50$_{\pm \phantom{0}.46}$ &          85.18$_{\pm .42}$ & - & - & - \\
        4 & \textbf{98.65}$_{\pm .04}$ &   \textbf{94.54}$_{\pm .25}$ &          95.94$_{\pm .06}$ & \textbf{84.62}$_{\pm \phantom{0}.86}$ & \textbf{85.23}$_{\pm .39}$ & - & - & - \\
        % 5 &          98.68$_{\pm .02}$ &            94.30$_{\pm .15}$ & \textbf{96.00}$_{\pm .07}$ &          84.44$_{\pm \phantom{0}.98}$ &          85.17$_{\pm .56}$ & - & - & - \\
        % 6 &          98.63$_{\pm .05}$ &            94.40$_{\pm .19}$ &          95.96$_{\pm .06}$ &          84.76$_{\pm \phantom{0}.52}$ & \textbf{85.91}$_{\pm .65}$ & - & - & - \\
        % 7 & \textbf{98.70}$_{\pm .06}$ &            94.60$_{\pm .10}$ &          95.92$_{\pm .07}$ & \textbf{85.28}$_{\pm \phantom{0}.32}$ &          85.77$_{\pm .61}$ & - & - & - \\
        \hline
        \multicolumn{5}{l}{Best results are in \textbf{bold}.}
    \end{tabular}
    \label{tab:abla-ngenerations}
\end{table*}

\end{comment}
